\documentclass[11pt,a4paper]{article}

\usepackage[T1]{fontenc}
\usepackage[utf8]{inputenc}
\usepackage[spanish,es-tabla]{babel}
\usepackage{lmodern}
\usepackage{microtype}
\usepackage{geometry}
\usepackage{booktabs}
\usepackage{tabularx}
\usepackage{array}
\usepackage{amsmath,amssymb}
\usepackage{setspace}
\usepackage{caption}
\usepackage{float}
\usepackage{enumitem}
\usepackage{xcolor}
\usepackage{hyperref}

\geometry{margin=2.0cm, top=1.7cm, bottom=1.8cm}
\setstretch{1.05}
\setlength{\parskip}{0.25em}
\setlength{\parindent}{0pt}
\setlist[itemize]{leftmargin=1.2em, itemsep=0.1em, topsep=0.15em}
\setlist[enumerate]{leftmargin=1.2em, itemsep=0.1em, topsep=0.15em}
\captionsetup{font=footnotesize, labelfont=bf, skip=3pt}
\hypersetup{colorlinks=true, linkcolor=black, urlcolor=blue!50!black, citecolor=black}

\newcommand{\pregunta}[1]{\noindent\textbf{Pregunta.} #1\par}
\newcommand{\resultado}[1]{\noindent\textbf{Resultado.} #1\par}
\newcommand{\interpretacion}[1]{\noindent\textbf{Interpretación.} #1\par}
\newcommand{\verif}[1]{\noindent\small\textit{Verificación y límites.} #1\par}

\title{\bfseries Problemas de transporte y flujo en redes\\
\large Balance, transbordo, faltantes penalizados y flujo de costo mínimo --- SaludRed / AyudaFederal}
\author{Lorenzo Marquesini\\[0.25em]
\small Investigación Operativa --- Clase 05}
\date{}

\begin{document}
\vspace{-0.6em}
\maketitle
\vspace{-1.4em}

\begin{abstract}
\noindent
Cuatro modelos de programación lineal continua resueltos con OR-Tools (GLOP): transporte balanceado desde dos plantas a tres centros; transbordo por un hub con conservación de flujo; transporte con demanda excedente y variables de faltante penalizadas; y flujo de costo mínimo en una red con fuente y sumidero. Se documentan formulaciones, restricciones de balance, soluciones óptimas y verificaciones de factibilidad.
\end{abstract}

\noindent\textit{Unidades:} lotes y centenas de unidades monetarias (costos de transporte y manipulación). \textit{Implementación:} \texttt{clase 5/main.ipynb}.

%==============================================================================
\section{C05E1 --- Transporte balanceado}
%==============================================================================

\subsection*{Situación}
SaludRed distribuye kits médicos desde dos plantas (P0, P1) hacia tres centros regionales (D0, D1, D2). P0 despacha hasta 35 lotes; P1 hasta 25. Los centros requieren exactamente 20, 25 y 15 lotes. Oferta total (60) coincide con demanda total (60): no hay inventario ni faltante.

\pregunta{¿Cuántos lotes debe enviar cada planta a cada centro para cubrir exactamente las tres demandas al menor costo semanal?}

\subsection*{Modelo}
Variables $x_{ij}\ge 0$: lotes desde planta $i$ hacia centro $j$ (continuas). Matriz de costos por lote:
\[
C = \begin{pmatrix}
4 & 6 & 8 \\
5 & 4 & 3
\end{pmatrix}
\quad\text{(filas P0, P1; columnas D0, D1, D2).}
\]
\begin{align}
\min \; Z &= \sum_{i,j} c_{ij}\, x_{ij},\\
x_{p0,d0} + x_{p0,d1} + x_{p0,d2} &= 35 && \text{(oferta P0)},\\
x_{p1,d0} + x_{p1,d1} + x_{p1,d2} &= 25 && \text{(oferta P1)},\\
x_{p0,d0} + x_{p1,d0} &= 20 && \text{(demanda D0)},\\
x_{p0,d1} + x_{p1,d1} &= 25 && \text{(demanda D1)},\\
x_{p0,d2} + x_{p1,d2} &= 15 && \text{(demanda D2)}.
\end{align}

\begin{table}[H]
\centering
\small
\caption{Arcos, costos y límites de oferta/demanda (C05E1).}
\label{tab:c05e1}
\begin{tabular}{@{}lcccc@{}}
\toprule
Arco & Costo & Oferta & Demanda \\
\midrule
p0-d0 & 4 & 35 & 20 \\
p0-d1 & 6 & 35 & 25 \\
p0-d2 & 8 & 35 & 15 \\
p1-d0 & 5 & 25 & 20 \\
p1-d1 & 4 & 25 & 25 \\
p1-d2 & 3 & 25 & 15 \\
\bottomrule
\end{tabular}
\end{table}

\subsection*{Resultado}
Estado: \textbf{óptimo} (viable). Solución:
\begin{table}[H]
\centering
\small
\caption{Flujos óptimos C05E1.}
\begin{tabular}{@{}lc@{}}
\toprule
Arco & Flujo \\
\midrule
p0-d0 & 20{,}00 \\
p0-d1 & 15{,}00 \\
p0-d2 & 0{,}00 \\
p1-d0 & 0{,}00 \\
p1-d1 & 10{,}00 \\
p1-d2 & 15{,}00 \\
\midrule
\multicolumn{2}{l}{Costo mínimo total = 255{,}0} \\
\bottomrule
\end{tabular}
\end{table}

Reconstrucción: $20{\cdot}4 + 15{\cdot}6 + 10{\cdot}4 + 15{\cdot}3 = 80 + 90 + 40 + 45 = 255$.

\interpretacion{P0 satura su oferta (35) enviando por rutas baratas hacia D0 y D1; P1 cubre el resto de D1 y toda D2 por la ruta más económica (p1-d2, costo 3). No se usa p1-d0 (costo 5) ni p0-d2 (costo 8).}

\verif{Ofertas: $20+15+0=35$; $0+10+15=25$. Demandas: $20+0=20$; $15+10=25$; $0+15=15$. Balance global: $60=60$.}

%==============================================================================
\section{C05E2 --- Transbordo por un hub}
%==============================================================================

\subsection*{Situación}
AyudaFederal traslada 50 lotes desde depósitos Oeste (O, 30 lotes) y Norte (N, 20) hacia regiones Centro (C) y Sur (S, 25 cada una), con posible consolidación en hub H. H no produce ni consume: flujo entrante = flujo saliente. Seis rutas dirigidas con costo y capacidad por arco.

\pregunta{¿Qué flujo debe circular por cada una de las seis rutas para abastecer C y S al menor costo sin violar capacidades ni crear unidades en H?}

\subsection*{Modelo}
Convención de balance: \textbf{salida $-$ entrada}. Nodos transitorios en 0; fuentes positivas; demandas negativas.
\begin{align}
\min \; Z &= 2\,x_{O\!-\!H} + 3\,x_{N\!-\!H} + 5\,x_{O\!-\!C} + 4\,x_{N\!-\!S} + x_{H\!-\!C} + 2\,x_{H\!-\!S},\\
0 &\le x_{arco} \le \text{capacidad}_{arco},\\
x_{O\!-\!H} + x_{N\!-\!H} &= x_{H\!-\!C} + x_{H\!-\!S} && \text{(hub H)},\\
x_{O\!-\!H} + x_{O\!-\!C} &= 30 && \text{(O)},\\
x_{N\!-\!H} + x_{N\!-\!S} &= 20 && \text{(N)},\\
x_{O\!-\!C} + x_{H\!-\!C} &= 25 && \text{(C)},\\
x_{N\!-\!S} + x_{H\!-\!S} &= 25 && \text{(S)}.
\end{align}

\begin{table}[H]
\centering
\small
\caption{Rutas, costos y capacidades (C05E2).}
\begin{tabular}{@{}lcc@{}}
\toprule
Arco & Costo & Capacidad \\
\midrule
O-H & 2 & 25 \\
N-H & 3 & 20 \\
O-C & 5 & 20 \\
N-S & 4 & 20 \\
H-C & 1 & 25 \\
H-S & 2 & 25 \\
\bottomrule
\end{tabular}
\end{table}

\subsection*{Resultado}
Estado: \textbf{óptimo}. Solución:
\begin{table}[H]
\centering
\small
\caption{Flujos óptimos C05E2.}
\begin{tabular}{@{}lc@{}}
\toprule
Arco & Flujo \\
\midrule
O-H & 25{,}00 \\
N-H & 0{,}00 \\
O-C & 5{,}00 \\
N-S & 20{,}00 \\
H-C & 20{,}00 \\
H-S & 5{,}00 \\
\midrule
\multicolumn{2}{l}{Costo mínimo total = 185{,}0} \\
\bottomrule
\end{tabular}
\end{table}

Reconstrucción: $25{\cdot}2 + 5{\cdot}5 + 20{\cdot}4 + 20{\cdot}1 + 5{\cdot}2 = 50 + 25 + 80 + 20 + 10 = 185$.

\interpretacion{O envía casi todo por el hub (ruta barata O-H, costo 2) y complementa C directamente (5 lotes). N va directo a S (N-S al máximo, 20) sin usar el hub. H transborda 25 entrantes hacia C (20) y S (5). La ruta H-C (costo 1) concentra el flujo consolidado hacia C.}

\verif{Hub: $25+0=20+5=25$. O: $25+5=30$. N: $0+20=20$. C: $5+20=25$. S: $20+5=25$. Todas las capacidades respetadas.}

%==============================================================================
\section{C05E3 --- Demanda no atendida con penalización}
%==============================================================================

\subsection*{Situación}
Contingencia en la misma red de C05E1: ofertas iguales (35 y 25), pero D1 aumenta de 25 a 40 lotes. Demanda total 75 vs.\ oferta 60. Se introducen variables $r_j\ge 0$ (lotes no atendidos en D0, D1, D2). Penalización común $p$ por lote faltante; el enunciado propone ensayar $p\in\{10,50,200\}$. El notebook resuelve $p=10$.

\pregunta{¿Cómo distribuir los 60 lotes disponibles y en qué centros registrar faltantes para minimizar costo logístico más penalización?}

\subsection*{Modelo}
\begin{align}
\min \; Z &= \sum_{i,j} c_{ij}\, x_{ij} + p\,(r_{d0}+r_{d1}+r_{d2}),\\
x_{p0,d0}+x_{p0,d1}+x_{p0,d2} &\le 35,\\
x_{p1,d0}+x_{p1,d1}+x_{p1,d2} &\le 25,\\
x_{p0,d0}+x_{p1,d0}+r_{d0} &= 20,\\
x_{p0,d1}+x_{p1,d1}+r_{d1} &= 40,\\
x_{p0,d2}+x_{p1,d2}+r_{d2} &= 15,\\
0 &\le r_{d0}\le 20,\quad 0\le r_{d1}\le 40,\quad 0\le r_{d2}\le 15.
\end{align}
Ofertas con $\le$ (no obligatorias); demandas con igualdad incluyendo faltante.

\subsection*{Resultado ($p=10$)}
Estado: \textbf{óptimo}.
\begin{table}[H]
\centering
\small
\caption{Solución C05E3 con penalización $p=10$.}
\begin{tabular}{@{}lc@{}}
\toprule
Concepto & Valor \\
\midrule
p0-d0 & 20{,}00 \\
p0-d1 & 15{,}00 \\
p0-d2 & 0{,}00 \\
p1-d0 & 0{,}00 \\
p1-d1 & 10{,}00 \\
p1-d2 & 15{,}00 \\
\midrule
$r_{d0}$ & 0{,}00 \\
$r_{d1}$ & 15{,}00 \\
$r_{d2}$ & 0{,}00 \\
\midrule
Costo logístico & 255{,}0 \\
Penalización ($10\times 15$) & 150{,}0 \\
Costo total & 405{,}0 \\
\bottomrule
\end{tabular}
\end{table}

\interpretacion{Con penalización baja ($p=10$), el modelo concentra el faltante en D1 (15 lotes): es más barato enviar por rutas económicas y pagar penalización que usar rutas caras para cubrir todo. D0 y D2 quedan totalmente servidos. El costo logístico (255) coincide con el óptimo balanceado de C05E1 sobre los 60 lotes disponibles.}

\verif{Ofertas: P0 usa 35; P1 usa 25 (saturación total). Demanda efectiva: $20+25+15=60$ enviados $+$ $15$ faltantes $=75$. Para $p=50$ y $p=200$ conviene reoptimizar: a mayor penalización, el faltante tendería a reducirse si existen rutas viables.}

%==============================================================================
\section{C05E4 --- Flujo de costo mínimo (fuente y sumidero)}
%==============================================================================

\subsection*{Situación}
Red con demanda 10 unidades desde fuente $s$ hasta sumidero $t$. Aristas: $(s,A)$, $(s,B)$, $(A,B)$, $(A,t)$, $(B,t)$ con costos $2,4,1,5,2$ y capacidades $6,8,3,6,8$.

\pregunta{¿Qué flujo por arista minimiza el costo total respetando balances nodales y capacidades?}

\subsection*{Balances nodales}
Convención: \textbf{flujo saliente $-$ flujo entrante}. Solo $s$ genera flujo neto ($+10$); solo $t$ lo absorbe ($-10$); nodos intermedios en balance cero.
\begin{align}
\text{nodo } s:&\quad x_{sA} + x_{sB} = 10,\\
\text{nodo } A:&\quad -x_{sA} + x_{AB} + x_{At} = 0,\\
\text{nodo } B:&\quad -x_{sB} - x_{AB} + x_{Bt} = 0,\\
\text{nodo } t:&\quad -x_{At} - x_{Bt} = -10.
\end{align}
$s$ es la única fuente porque debe inyectar 10 unidades netas hacia la red; $t$ es el único sumidero porque debe recibir esas 10 unidades. $A$ y $B$ solo redistribuyen: todo lo que entra debe salir.

\subsection*{Modelo}
\begin{equation}
\min \; Z = 2\,x_{sA} + 4\,x_{sB} + x_{AB} + 5\,x_{At} + 2\,x_{Bt},
\end{equation}
con $0\le x_{arista}\le \text{capacidad}_{arista}$.

\begin{table}[H]
\centering
\small
\caption{Aristas, capacidades y costos (C05E4).}
\begin{tabular}{@{}lcc@{}}
\toprule
Arista & Capacidad & Costo \\
\midrule
s-A & 6 & 2 \\
s-B & 8 & 4 \\
A-B & 3 & 1 \\
A-t & 6 & 5 \\
B-t & 8 & 2 \\
\bottomrule
\end{tabular}
\end{table}

\subsection*{Resultado}
Estado: \textbf{óptimo}.
\begin{table}[H]
\centering
\small
\caption{Flujos óptimos C05E4.}
\begin{tabular}{@{}lc@{}}
\toprule
Arista & Flujo \\
\midrule
s-A & 5{,}00 \\
s-B & 5{,}00 \\
A-B & 3{,}00 \\
A-t & 2{,}00 \\
B-t & 8{,}00 \\
\midrule
\multicolumn{2}{l}{Costo mínimo total = 59{,}0} \\
\bottomrule
\end{tabular}
\end{table}

Reconstrucción: $5{\cdot}2 + 5{\cdot}4 + 3{\cdot}1 + 2{\cdot}5 + 8{\cdot}2 = 10 + 20 + 3 + 10 + 16 = 59$.

\interpretacion{Se usa el arco barato $A$-$B$ (costo 1) al máximo (3). $s$ reparte 5+5 entre $A$ y $B$. Desde $A$, 2 van directo a $t$ y 3 pasan por $B$ vía $A$-$B$; desde $B$, $3+5=8$ salen hacia $t$ por la ruta económica $B$-$t$.}

\verif{Balance $s$: $5+5=10$. $A$: $-5+3+2=0$. $B$: $-5-3+8=0$. $t$: $-2-8=-10$. Capacidades: $s$-$A$ usa 5/6; $A$-$B$ satura 3/3; $B$-$t$ satura 8/8.}

%==============================================================================
\section{Síntesis}
%==============================================================================

\begin{table}[H]
\centering
\small
\caption{Resumen de los cuatro ejercicios.}
\label{tab:sint}
\begin{tabularx}{\textwidth}{@{}lXXX@{}}
\toprule
Ejercicio & Problema & Óptimo & Idea clave \\
\midrule
C05E1 & Transporte balanceado & $Z=255$ & Igualdades en oferta y demanda \\
C05E2 & Transbordo hub & $Z=185$ & Conservación en H; rutas mixtas \\
C05E3 & Faltante penalizado ($p=10$) & $Z=405$ & $r_{d1}=15$; oferta $\le$ \\
C05E4 & Flujo $s\to t$ & $Z=59$ & Balance nodal; $A$-$B$ al tope \\
\bottomrule
\end{tabularx}
\end{table}

\subsection*{Conclusiones}
\begin{enumerate}
\item En transporte balanceado, las restricciones de planta son igualdades (toda la oferta se despacha) y las de destino garantizan servicio exacto.
\item El hub impone conservación de flujo; explicitar la convención de signos evita errores al auditar nodos transitorios.
\item Las variables de faltante convierten un problema infactible en factible, separando costo logístico de penalización por nivel de servicio.
\item En flujo de costo mínimo, solo la fuente y el sumidero tienen balance distinto de cero; los nodos intermedios conservan flujo.
\end{enumerate}

\subsection*{Supuestos}
\begin{itemize}
\item Variables de flujo continuas (lotes agregados divisibles).
\item Costos lineales por unidad transportada; penalización lineal homogénea en C05E3.
\item Capacidades de arco independientes; sin costos fijos ni indivisibilidades.
\item C05E3: penalización idéntica en todos los centros (no refleja prioridades sanitarias diferenciadas).
\end{itemize}

\vspace{0.5em}
\noindent\small\textit{Solver:} OR-Tools GLOP (\texttt{pywraplp}). \textit{Fuente:} \texttt{clase 5/main.ipynb}.

\end{document}
